\chapter{Motor de reglas: Kids on Bikes}

\section{Filosofía del sistema}

\emph{Kids on Bikes} es un juego de rol narrativo y accesible diseñado para recrear historias de aventuras juveniles, bicicletas y misterios sobrenaturales en pequeñas comunidades, en la línea de los clásicos cinematográficos de los años 80 y 90.

\section{Los seis atributos y la escala de dados}

En lugar de modificadores numéricos fijos, cada personaje asigna un tipo de dado poliédrico distinto a cada uno de sus seis atributos:
\[
\mathbf{d20,\quad d12,\quad d10,\quad d8,\quad d6,\quad d4}
\]
Un dado más grande representa mayor talento o entrenamiento natural en esa disciplina.

\begin{itemize}
  \item \lead{Astucia (Brains):} Deducción lógica, memoria, cultura general, arreglar aparatos electrónicos y entender planos o acertijos.
  \item \lead{Músculo (Brawn):} Fuerza física pura, empujar objetos pesados, nadar contra corriente o forzar una reja.
  \item \lead{Combate (Fight):} Capacidad para defenderse en una pelea cuerpo a cuerpo, lanzar piedras con precisión o forcejear.
  \item \lead{Reflejos (Flight):} Velocidad de reacción, esquivar proyectiles, mantener el equilibrio y pedalear a toda velocidad en bicicleta.
  \item \lead{Vigor (Grit):} Dureza mental, agallas para no salir huyendo presa del pánico, sangre fría y resistencia al cansancio.
  \item \lead{Labia (Charm):} Carisma, empatía, convencer a los adultos, mentir a la policía o calmar a un compañero en crisis.
\end{itemize}

\section{Mecánica de resolución y dados explosivos}

Cuando un personaje intenta una acción cuyo desenlace es incierto o peligroso, el Director fija una \lead{Dificultad} numérica y el jugador lanza el dado correspondiente a su atributo:

\begin{CajaRegla}{Escala de dificultad estándar}
\begin{itemize}
  \item \textbf{1--3:} Tarea cotidiana que requiere algo de atención.
  \item \textbf{4--6:} Tarea estándar bajo presión moderada.
  \item \textbf{7--9:} Desafío difícil; requiere experiencia o agallas.
  \item \textbf{10--14:} Muy difícil; reservado a especialistas o grandes golpes de suerte.
  \item \textbf{15--19:} Casi imposible; desafía las leyes ordinarias.
  \item \textbf{20+:} Hazaña legendaria; proeza cinematográfica.
\end{itemize}
\end{CajaRegla}

\begin{CajaRegla}{Dados explosivos (\emph{Exploding Dice})}
Si un jugador obtiene el valor numérico más alto posible en su dado (por ejemplo, un $6$ en 1d6, o un $12$ en 1d12), \textbf{el dado explota}: el jugador anota el resultado y vuelve a lanzar el mismo dado, sumando el nuevo valor al total. ¡Este proceso se repite tantas veces como continúe saliendo el valor máximo!
\end{CajaRegla}

\section{Fichas de Adversidad (\emph{Adversity Tokens})}

En \emph{Kids on Bikes}, el fracaso impulsa la historia hacia adelante:
\begin{itemize}
  \item Cada vez que un personaje intenta una tirada y no alcanza la Dificultad establecida, \lead{gana 1 Ficha de Adversidad}.
  \item Un jugador puede gastar 1 Ficha de Adversidad antes o después de una tirada para sumar un $+1$ al resultado final.
  \item Varios jugadores pueden unir sus Fichas de Adversidad para ayudar a un compañero en un momento crítico.
  \item \textbf{Uso en el Personaje Potenciado:} Las Fichas de Adversidad son el recurso principal para pagar el uso de los poderes psíquicos de Alex Etxebarria.
\end{itemize}

\section{Vehículos y bicicletas}

Las bicicletas son el transporte por excelencia de los jóvenes de Sestao. Cada jugador describe su bicicleta (modelo, color, trucos o parches en las ruedas) y anota su ventaja:
\begin{itemize}
  \item \lead{BH California / Orbea BMX:} Otorga un $+1$ a las tiradas de \emph{Reflejos} al saltar bordillos, descender terraplenes o hacer maniobras evasivas en callejones estrechos.
  \item \lead{Bicicleta de montaña / Paseo con transportín:} Permite llevar a un segundo pasajero o transportar herramientas y equipo pesado sin penalización.
\end{itemize}
